\documentclass[10pt]{article}
\usepackage[margin=0.9in]{geometry}
\usepackage{booktabs,hyperref,amsmath}
\newcommand{\todo}[1]{\textbf{[TODO: #1]}}
\newcommand{\unread}{$^{\dagger}$}
\title{TripleLight: an Embedded Analytical RDF Database}
\author{TripleLight authors}
\date{Draft, 2026-11-22 (target: VLDB 2027, February 1 cycle). Section~4 awaits the dedicated-machine measurements.}
\begin{document}
\maketitle

\noindent\textbf{Status of this draft.} Every number cites its evidence file (\texttt{docs/...} in the repository).
Papers marked~\unread\ were not read in full; they are named from their abstracts or from a source we read and must be re-checked before
submission. Timing results are deliberately absent: they are measured once, on one dedicated machine, under a registered protocol
(\texttt{docs/eval-plan/PLAN.md}). Per the VLDB 2027 submission guidelines (checked by Jianlin): research review is single-blind, so author names and affiliations appear on page one; regular research papers may use 12 pages excluding references; the February~1 cycle needs an abstract by January~25 and the March~1 cycle one by February~25.

\begin{abstract}
Applications that need analytical queries over RDF data today run a database server or load the data into a general-purpose
relational engine. TripleLight is an embedded RDF database: an application opens a directory through a library call, and
SPARQL~1.1 queries run in the application's process, over dictionary-encoded IDs, on a sorted, compressed, disk-resident store.
\todo{one sentence on the measured position against embedded RDF stores and a SQL engine (C1a), and one against RDF servers (C1b), from Section~4 after the dedicated-machine run.}
We describe the architecture and the techniques it combines, credit each to its source, and examine whether the comparison itself is fair, since
engines differ in which triple-index orders they store.
\end{abstract}

\section{Introduction}
\paragraph{Analytical SPARQL inside the application.} Data scientists increasingly run analyses on one machine and are poorly
served by systems built for dedicated servers, a point Raasveldt and M\"uhleisen made for relational data when they proposed
embedded analytical systems: tight integration with the analysis tool, efficient transfer of results into it, resilience on user
hardware, and cooperation with the host application for resources~\cite{raasveldt2020}. RDF has the same gap, and it is only partly filled. The scalable RDF stores we compare against (Virtuoso, QLever) are servers that an application reaches over a protocol.
Recent work moves toward the embedded setting from three sides. A community extension for DuckDB runs SPARQL~1.1 queries (a subset) over DuckDB tables through R2RML mappings~\cite{duckdbrdf}.
RDF Fusion builds a SPARQL engine on Apache DataFusion and Arrow; it keeps its quad indexes (GSPO, GPOS, GOPS) in memory, has no on-disk storage yet, and was evaluated on BSBM data of about 9 million and 1.8 million triples~\cite{schwarzinger2025}.
MillenniumDB is a persistent, open-source graph database that combines B+-tree indexes with worst-case-optimal joins, and lists SPARQL~1.1 among its supported models~\cite{vrgoc2021,mdbreadme}; it runs as a database server, and its published evaluation is on Wikidata with its own query language.
To our knowledge none of these is an embedded, disk-resident, SPARQL-native engine measured at hundreds of millions of triples, which is the position TripleLight takes. \todo{verify against each system's current release before submission, and add the measured comparison from Section~4.}

\paragraph{What we built and where it stands.} TripleLight is an embedded RDF database in Rust, with bindings for Python, Node, WebAssembly, C, Go and Java.
It answers SPARQL~1.1 queries and updates in process, on IDs, and returns results as SPARQL JSON, owned rows or Arrow C streams
(\texttt{DESIGN.md}). We measure it in two settings that are never mixed: embedded, through each language's in-process API, against embedded
RDF stores and a SQL engine (C1a); and as an HTTP SPARQL server, against QLever, Virtuoso, Jena TDB2/Fuseki, Oxigraph and RDF4J (C1b).
\todo{state the measured position for C1a and C1b, with the dedicated-machine numbers and the answer gate.}
The techniques that produce it are not new, and Section~3 credits each: four of the six triple orders with exact derivation of the other two,
count answers from index metadata, leapfrog triejoin with characteristic-set planning, and an external-sort build pipeline.
What we contribute is the combination in an embedded, disk-resident engine, the measurements, and the evidence package behind them.

\paragraph{Is the comparison fair?} RDF engines differ in the triple orders they store (RDF-3X and Hexastore keep six; QLever supports a two-order mode; RDF4J's native store is configurable).
A ranking of engines can therefore mix executor quality with index choice. Section~4.5 asks whether changing the registered layouts of three engines moves them in
a ranking of seven HTTP engines, and registers in advance what would count as a negative answer.

\section{Architecture}
TripleLight is a Cargo workspace of focused crates (\texttt{DESIGN.md}): RDF terms and parsers, IRI handling, a core of term IDs and storage traits,
a SPARQL parser and algebra, a planner with statistics, an executor, a storage layer, and the embedded API with an optional HTTP server.
No Oxigraph crate is a runtime dependency; the parsers and term types were written to match \texttt{oxrdf}, \texttt{oxiri} and \texttt{oxttl} and are tested
differentially against them.

\paragraph{Terms and IDs.} Every term is a 64-bit ID whose top four bits give its kind. Within a kind, ID order equals term-text order, so operators compare, join and group
IDs and decode a term only when a filter, a string function or the result needs its text. On disk one front-coded dictionary file in ID order serves lookups in both directions;
this one copy shrank the BSBM-100M store from 12,533.1\,MB to 7,936.3\,MB (\texttt{docs/unified-dictionary-100m-v1}).

\paragraph{Storage.} A base snapshot is immutable: sorted triples in compressed, checksummed blocks, read on demand through decoded-block caches that share one memory budget.
New bases keep SPO, PSO, POS and OSP; SOP and OPS are served from them with the same rows in the same order. Updates go to a write-ahead log with group commit and form an
overlay of inserts and deletes over the base; readers pin a snapshot, so a query never sees part of an update, and compaction writes a new base by external sort and
switches to it atomically. A file lock allows one process per store directory.
Commits append to the log and fsync before returning, and concurrent commits share one fsync. Compaction publishes the new base by fsyncing it, atomically renaming \texttt{CURRENT}, fsyncing the directory, and only then truncating the log. We test the crash states between these steps (\texttt{compaction\_crash\_states\_recover}, torn log tails); we did not run filesystem fault injection. TripleLight targets read-mostly analytics: updates land in an in-memory delta over an immutable base, compaction rewrites the base ($O$(base)), and we do not evaluate sustained update throughput (statements supplied by the storage owner, task \#100).

\paragraph{Query processing.} A BGP is ordered greedily by exact index-prefix counts, and eligible patterns use a cost-based order from characteristic-set statistics; cyclic patterns use leapfrog triejoin.
Execution runs on IDs with index nested-loop joins, filter pushdown, hash aggregation, top-$k$ for \texttt{ORDER BY \dots\ LIMIT}, and count shortcuts. Statistics are persisted next to the base:
with a saved file they were ready 0.74 to 0.91\,s after open, against 20.5 to 25.8\,s when rebuilt by a scan (\texttt{docs/startup36-persisted-stats}; laptop, single session, preliminary: to be replaced by dedicated-machine numbers).
Execution is intra-query parallel through Rayon: morsel-parallel BGP extension, parallel counting at the last join level, a parallel first level for leapfrog, and parallel result decoding and serialization; aggregation and the other operators run serially (code audit in task \#100; a thread-scaling run is planned on the dedicated machine). There is one writer and one process per store.
Correctness is gated on answers matching Apache Jena under each benchmark's declared checks and on at least 323 passing W3C SPARQL~1.1 query tests in CI (\texttt{DESIGN.md}).

\section{Techniques and Their Sources}
\subsection{Four orders with exact derivation}
RDF-3X stores all six permutations of subject, predicate and object, compressed by deltas, and adds count-aggregated indexes~\cite{neumann2008}; Hexastore indexes the data in all six orders, at the price of a worst-case five-fold increase in index space~\cite{weiss2008}.
QLever's authors report that for typical semantic queries two permutations (PSO and POS) suffice, at the loss of patterns with a variable predicate, and describe Virtuoso's default indexes as corresponding to the same two, with more permutations built on demand~\cite{bast2017}.
TripleLight keeps four (SPO, PSO, POS, OSP) and derives SOP and OPS exactly, so every access order remains available. The default store is 36\% smaller on LUBM-1000 (4.97 to 3.16\,GB) and 17\% smaller
on BSBM-100M (7.94 to 6.56\,GB) than the six-order store (\texttt{docs/build-compare-orders-v1}). Across 255 benchmark queries, the trace recorded no regroup scans or memo misses and 18 direct remaps, all on LUBM (\texttt{docs/ablation-qual-v1}, \todo{cite after PR \#146 merges}); any latency effect is left to the timed four-against-six comparison in Section~4.4. How often real query logs reach a derived order is open (Section~5).
The ring of Arroyuelo et al.\ replaces the six orders needed for worst-case-optimal joins by one compact structure that adds only sublinear space; their implementation operates in main memory on Wikidata, and they note that its random accesses make migration to disk difficult~\cite{arroyuelo2021}. Our setting differs: disk-resident, ordinary sorted blocks, and four stored orders.

\subsection{Counting from index metadata and ID-level aggregation}
Aggregation runs on IDs and decodes only the output. For supported counts the executor reads exact range counts from the index instead of enumerating matches:
a single-pattern \texttt{COUNT} is a prefix count, and a group-by on a permutation prefix counts each run (\texttt{docs/count-pushdown}). RDF-3X precomputes such counts as separate aggregated indexes~\cite{neumann2008};
ours use the exact prefix metadata of the stored orders and need no extra index. Kùzu's factorized query processor represents intermediate relations of many-to-many joins as Cartesian products instead of flat tuples~\cite{feng2023}, which is closely related to counting without enumerating. The effect is measured as an ablation (Section~4.6), not claimed as a contribution.

\subsection{Leapfrog triejoin with characteristic-set planning}
Cyclic patterns use leapfrog triejoin~\cite{veldhuizen2014}, a worst-case-optimal join that requires its inputs as sorted iterators, so each variable order needs an index in the matching permutation. The planner estimates cardinalities with characteristic sets and predicate counts (Neumann and Moerkotte\unread), orders joins by subset dynamic programming, and picks a variable order for trie joins;
it replaces the greedy order only when statistics exist, the pattern has two to seven triples and the estimated greedy work is at least 100,000 (\texttt{docs/optimizer.md}).
Its costs are heuristic work units, not calibrated against elapsed time (\texttt{docs/optimizer.md}).

\subsection{Build pipeline}
A disk store is built by external sort: parallel shard parsing, a term sort, ID assignment and the order builds. On LUBM-1000 a four-order build took 654.8\,s on the development laptop (preliminary, single session), of which shard parsing and sorting took 38.7\%, the dictionary-side stages 35\% and the order builds 26.4\%
(\texttt{docs/build-profile-v1}). Fixed-width sort rows were 4--5$\times$ faster than variable-length rows in a laptop microbenchmark (same file; preliminary). Terms repeat heavily: 7.9\% are distinct in LUBM-1000 and BSBM-100M and 3.2\% in WatDiv-100 over the whole file
(\texttt{docs/parse-phase-v1}), which makes deduplication before sorting attractive; QLever's builder already deduplicates before sorting, so we treat this as engineering that follows an existing design.

\subsection{Memory budget and cooperation}
The decoded-block caches share one budget. Cache budgets of 2\,GB, 1\,GB and 512\,MB made warm LUBM-1000 queries 3.6, 4.0 and 4.6$\times$ slower than the default on the development laptop (\texttt{docs/cache-budget-sweep-v1}; preliminary, to be replaced by dedicated-machine numbers). Cooperating with a host application's resources is one of the embedded requirements~\cite{raasveldt2020};
we do not yet claim more than a configurable budget. \todo{cgroup-aware default, once its PR is merged and measured.}

\section{Evaluation}
\subsection{Setup and answer gates}
\todo{Systems, versions, dedicated machine, protocol, answer classes and timeout: from \texttt{docs/eval-plan/PLAN.md}.}\\
Memory rule (decided in task \#100): every system runs in a cgroup with the same \texttt{memory.max}; inside it each system's own memory setting follows its vendor's guidance, stated per system (TripleLight: half of the limit for its block cache).\\
Commercial engines (GraphDB, Stardog, AllegroGraph) are excluded: we do not accept licence terms on publishing benchmark results on the owner's behalf.
\subsection{C1a: embedded}\todo{Results after the dedicated-machine run.}
\subsection{C1b: HTTP}\todo{Results after the dedicated-machine run.}
\subsection{Build and storage}\todo{Build time, scratch and store size on the dedicated machine; size numbers from Section~3.1 may be repeated.}
\subsection{Is the comparison fair?}
Registered in the plan (\texttt{docs/eval-plan/PLAN.md}, amendment A4, merged 2026-11-22), before any measurement. The HTTP cohort is TripleLight, QLever, Virtuoso, Jena TDB2, Oxigraph, RDF4J and MillenniumDB; each must pass its store receipt and answer gate.
The moving engines are TripleLight (four and six orders), QLever (default and two-order) and RDF4J (default \texttt{spoc,posc} and aligned \texttt{spoc,posc,ospc}); the others are fixed-layout.
The statistic is Kendall's $\tau_b$ with average ranks, per suite and per independent block, on one frozen common set of engines and requests used in both rank vectors; the suite reducer is the median of the per-block values, with every block shown.
A suite needs at least four paired engines after unsupported cells are dropped from both vectors, otherwise it is unevaluable, and if no suite is evaluable the test is unevaluable. The claim fails if $\tau_b \ge 0.8$ on every evaluable suite.
Registered prediction (claims owner, 2026-11-22): $\tau_b < 0.8$ on BSBM BI and the WatDiv aggregates, and $\tau_b \ge 0.8$ on LUBM. If a prediction is wrong, the paper reports that.
\todo{results table: per suite, paired engines, per-block $\tau_b$, median, the ranks that change.}
\subsection{Ablations}\todo{Count shortcut, cache budget, result windowing: reported as ablations.}

\section{Limits and Related Work}
\paragraph{Limits.} Disk builds parse shards in parallel and overlap index builds (an average of 2.9 of 10 cores busy on LUBM-1000, \texttt{docs/build-profile-v1}) and need scratch disk of 1.7 to 2.1 times the input at the peak (\texttt{docs/build-compare-orders-v1}).
Results are materialized before output. One process per store directory. \texttt{LOAD}, \texttt{SERVICE}, \texttt{BNODE(string)} and blank-node graph names are unsupported (\texttt{DESIGN.md}).
The cost of the derived orders on real query logs is not measured. The benchmarks are synthetic and their structure is nearly fixed and unlike real RDF data~\cite{duan2011}.
\paragraph{Related work.} \emph{Embedded analytics.} DuckDB's requirements for embedded analytical systems~\cite{raasveldt2020} (the system itself is described in a SIGMOD 2019 demonstration paper~\cite{raasveldt2019}); Kùzu~\cite{feng2023} (a disk-based, embeddable graph database queried in Cypher, so outside the SPARQL comparison; \todo{confirm its current maintenance status before stating it}); pyoxigraph, Blazegraph and DuckDB as embedded comparators.
\emph{RDF stores.} RDF-3X~\cite{neumann2008}, Hexastore~\cite{weiss2008}, QLever~\cite{bast2017}, Virtuoso, Jena TDB2, Oxigraph, RDF4J.
\emph{Layout and benchmarks.} Sidirourgos et al.\ re-evaluated vertical partitioning on MonetDB/SQL and PostgreSQL and found that clustering (PSO against SPO) changes performance by factors of two to five~\cite{sidirourgos2008};
Duan et al.\ showed that benchmark data differ from real RDF data in structuredness~\cite{duan2011}.

\section*{Artifact}
\todo{Move to the official PVLDB Vol.~20 template (github.com/vldbproceedings/VLDB-Template): set the public archive URL with \textbackslash vldbavailabilityurl and keep \textbackslash vldbtopmatter after \textbackslash maketitle, so Artifact Availability is rendered in the front matter. Regular Research needs supplemental materials in a public archive or a stated reason for an exception, per Jianlin's reading of the guidelines.}
The code, the scripts and the data needed to reproduce every number are in the repository \todo{URL} and in the archived package \texttt{docs/eval-v1/} (raw receipts, driver output, dataset checksums and reproduction commands), produced in the dedicated-machine window. \todo{link both, and say which commit and which package version the paper uses; VLDB requires the link in the initial submission.}

\begin{thebibliography}{9}
\bibitem{raasveldt2020} M.~Raasveldt, H.~M\"uhleisen. Data management for data science: towards embedded analytics. \emph{CIDR}, 2020.
\bibitem{neumann2008} T.~Neumann, G.~Weikum. RDF-3X: a RISC-style engine for RDF. \emph{PVLDB} 1(1), 2008.
\bibitem{bast2017} H.~Bast, B.~Buchhold. QLever: a query engine for efficient SPARQL+text search. \emph{CIKM}, 2017.
\bibitem{sidirourgos2008} L.~Sidirourgos, R.~Goncalves, M.~Kersten, N.~Nes, S.~Manegold. Column-store support for RDF data management: not all swans are white. \emph{PVLDB} 1(2), 2008.
\bibitem{schwarzinger2025} T.~Schwarzinger, M.~Thoma, T.~Preindl, M.~Kj{\"a}er, V.~P.~Just, G.~Steindl. RDF Fusion: an extensible SPARQL engine for hybrid data models. \emph{IEEE Access} 13, 2025.
\bibitem{vrgoc2021} D.~Vrgo{\v c} et al. MillenniumDB: a persistent, open-source, graph database. arXiv:2111.01540, 2021.
\bibitem{mdbreadme} MillenniumDB project README, github.com/MillenniumDB/MillenniumDB (accessed 2026-11-22).
\bibitem{duckdbrdf} DuckDB community extension \texttt{rdf} (``Read and write RDF, query using SPARQL''), duckdb.org/community\_extensions/extensions/rdf (accessed 2026-11-22).
\bibitem{arroyuelo2021} D.~Arroyuelo, A.~Hogan, G.~Navarro, J.~L.~Reutter, J.~Rojas-Ledesma, A.~Soto. Worst-case optimal graph joins in almost no space. \emph{SIGMOD}, 2021.
\bibitem{veldhuizen2014} T.~L.~Veldhuizen. Leapfrog triejoin: a simple, worst-case optimal join algorithm. \emph{ICDT}, 2014 (arXiv:1210.0481).
\bibitem{weiss2008} C.~Weiss, P.~Karras, A.~Bernstein. Hexastore: sextuple indexing for Semantic Web data management. \emph{PVLDB} 1(1), 2008.
\bibitem{feng2023} X.~Feng, G.~Jin, Z.~Chen, C.~Liu, S.~Salihoglu. K\`uzu graph database management system. \emph{CIDR}, 2023.
\bibitem{raasveldt2019} M.~Raasveldt, H.~M\"uhleisen. DuckDB: an embeddable analytical database. \emph{SIGMOD} (demonstration), 2019.
\bibitem{duan2011} S.~Duan, A.~Kementsietsidis, K.~Srinivas, O.~Udrea. Apples and oranges: a comparison of RDF benchmarks and real RDF datasets. \emph{SIGMOD}, 2011.
\end{thebibliography}
\end{document}
